% Optional introduction paragraphs. Use selectively after the fresh proof reviews.
% Bibliography keys are defined in references.bib.

Stein 与 Strömberg 对一般对称凸体的中心极大算子建立了量级为
$n\log n$ 的弱型 $(1,1)$ 估计，其证明通过有限球族的选择及加权重叠控制完成
\cite[Section~2, pp.~260--263]{SteinStromberg1983}。
本文沿用这一有限选择框架，引入独立的选择阈值和缓冲膨胀参数，
并利用轴平行立方体的坐标结构估计其余增量的加权体积，
由此得到可按维数直接使用的显式上界公式。

Aldaz 对 $\ell^\infty$ 几何下的中心极大算子给出了关于测度的统一估计，
其常数为 $2^n$\cite[Corollary~3.5]{Aldaz2020Kissing}。
取 Lebesgue 测度便得到本文模型下的基准 $c_n\le 2^n$。
本文以这一有明确出处的界作为七维应用的比较基准。

% Optional background paragraph; omit if asymptotic dimension dependence is not discussed.
Aldaz 证明了中心轴平行立方体极大算子的最优弱型 $(1,1)$ 常数随维数趋于无穷
\cite{Aldaz2011Cubes}。
本文研究给定维数下可核验的显式上界。

% The seven-dimensional decimal belongs to the paper's proved application/certificate,
% not to any of the three external sources. This literature review did not compute it.
